\documentclass[11pt,a4paper]{article}
% Reuse the visual language of the first exercise session.
\usepackage[margin=18mm,headheight=24pt]{geometry}
\usepackage[T1]{fontenc}
\usepackage[english]{babel}
\usepackage{lmodern}
\usepackage{microtype}
\usepackage{amsmath,amssymb}
\usepackage{siunitx}
\usepackage{tikz}
\usetikzlibrary{arrows.meta,calc,angles,quotes}
\usepackage{enumitem}
\usepackage{xcolor}
\usepackage{fancyhdr}
\usepackage{lastpage}
\usepackage[hidelinks]{hyperref}

\definecolor{VectorGreen}{HTML}{174A38}
\definecolor{VectorAccent}{HTML}{216343}
\definecolor{VectorMint}{HTML}{E2F0E5}
\definecolor{VectorCoral}{HTML}{C45136}
\definecolor{VectorInk}{HTML}{203B31}
\definecolor{VectorMuted}{HTML}{5F7468}

\color{VectorInk}
\setlength{\parindent}{0pt}
\setlength{\parskip}{0.55em}
\setlist[enumerate]{leftmargin=*,itemsep=0.45em,topsep=0.35em}
\sisetup{per-mode=symbol}

\pagestyle{fancy}
\fancyhf{}
\lhead{\textcolor{VectorMuted}{Vectors · Exercise session}}
\rhead{\textcolor{VectorMuted}{Addition and subtraction}}
\cfoot{\textcolor{VectorMuted}{\thepage\ / \pageref*{LastPage}}}
\renewcommand{\headrulewidth}{0.4pt}
\renewcommand{\headrule}{\hbox to\headwidth{\color{VectorMint}\leaders\hrule height \headrulewidth\hfill}}

\newcommand{\sessiontitle}[2]{%
  {\color{VectorGreen}\fontsize{25}{29}\selectfont\bfseries #1\par}
  \vspace{0.2em}
  {\color{VectorMuted}\large #2\par}
  \vspace{0.55em}
  {\color{VectorAccent}\rule{\textwidth}{1.5pt}}
}

\newcommand{\problemheading}[2]{%
  {\color{VectorAccent}\small\bfseries\MakeUppercase{Problem #1}}\par
  {\color{VectorGreen}\LARGE\bfseries #2}\par
  \vspace{0.25em}
}

\newcommand{\prediction}[1]{%
  \begin{center}
  \colorbox{VectorMint}{\parbox{0.92\linewidth}{\textbf{Predict before calculating.} #1}}
  \end{center}
}

\newcommand{\checkpoint}[1]{%
  \begin{center}
  \fcolorbox{VectorAccent}{white}{\parbox{0.91\linewidth}{\textbf{Reasoning checkpoint.} #1}}
  \end{center}
}

\newcommand{\worklines}[1]{%
  \par\vspace{0.25em}
  {\color{VectorMuted}\small\textit{Working and explanation}}\par
  \foreach \n in {1,...,#1}{\vspace{0.54cm}\noindent\textcolor{VectorMint}{\rule{\linewidth}{0.35pt}}\par}
}

\newcommand{\answerbox}[1]{%
  \begin{center}
  \colorbox{VectorMint}{\parbox{0.92\linewidth}{\textbf{Answer.} #1}}
  \end{center}
}

\lhead{\textcolor{VectorMuted}{Motion · Exercise session 2}}
\rhead{\textcolor{VectorMuted}{Constant acceleration}}

\rhead{\textcolor{VectorMuted}{Worked solutions}}

\begin{document}
\sessiontitle{Solutions: Motion hides in the timing}{Instructor copy · Exercise session 2}

\problemheading{1}{Three snapshots, one missing story}

\textbf{a) Reconstruct the motion.} With $\vec r_0=\vec0$, snapshots B and C give, in SI units,  
\[
\vec v_0+\tfrac12\vec a=(6,2),\qquad
3\vec v_0+\tfrac92\vec a=(0,12).
\]
Subtract three times the first equation from the second:
$3\vec a=(-18,6)$. Hence
\[
\boxed{\vec v_0=(9,1)\ \si{m/s},\qquad \vec a=(-6,2)\ \si{m/s^2}.}
\]
For use below, with $t$ in seconds,
\[
x=9t-3t^2,\quad y=t+t^2,\qquad v_x=9-6t,\quad v_y=1+2t.
\]

\textbf{b) A turn is not a stop.} The rightmost position occurs where $v_x=0$:
\[
\boxed{t=\SI{1.5}{s}.}
\]
Since $a_x<0$, the robot moves right before this time and left after it. At the turn,
\[
\vec v=(0,4)\ \si{m/s}.
\]
It is still moving in the positive $y$ direction, so it has not stopped.

\textbf{c) Average velocity.} Divide the displacement from A to C by the full elapsed time:
\[
\boxed{\bar{\vec v}=\frac{(0,12)-(0,0)}{3}=(0,4)\ \si{m/s}.}
\]
\checkpoint{The snapshots are separated by 1 s and 2 s. They are not equally spaced in time.}

\clearpage
\problemheading{2}{A throw through a moving window}

\textbf{a) Horizontal meeting.} The screen is at $x_s=8+2t$ and the ball at $x_b=6t$, with SI units understood. Thus
\[
6t=8+2t\quad\Longrightarrow\quad \boxed{t_c=\SI{2}{s}.}
\]
Both are then at $x=\SI{12}{m}$.

\textbf{b) Clear the edges.} The ball's height is $y=ut-5t^2$, so at the crossing $y_c=2u-20$. The window requires
\[
3<2u-20<5\quad\Longrightarrow\quad
\boxed{\SI{11.5}{m/s}<u<\SI{12.5}{m/s}.}
\]
The inequalities are strict because touching an edge does not count. The ball returns to ground at $t_g=2u/g$, between \SI{2.3}{s} and \SI{2.5}{s}, so every permitted launch is still airborne at the crossing.

\textbf{c) Centre the throw.} The centre is at $y=\SI{4}{m}$. Set $2u-20=4$:
\[
\boxed{u=\SI{12}{m/s}.}
\]
At the window, the vertical velocity is
\[
v_y=u-gt_c=12-10\cdot2=\SI{-8}{m/s}.
\]
The ball is \textbf{falling} as it passes through the centre.

\checkpoint{Match horizontal positions first. Then use that same time to check the height.}

\clearpage
\problemheading{3}{The slowest possible interception}
{\color{VectorCoral}\bfseries Optional challenge}

\textbf{a) Components and deadline.} Use $s$ uphill and $n$ normally away from the ramp. The target and ball obey
\[
s_T=L-\tfrac12g\sin\alpha\,t^2,\qquad n_T=0,
\]
\[
s_B=v_{0s}t-\tfrac12g\sin\alpha\,t^2,\qquad
n_B=v_{0n}t-\tfrac12g\cos\alpha\,t^2.
\]
Equating their positions at the meeting gives
\[
\boxed{v_{0s}=\frac Lt,\qquad v_{0n}=\frac{g\cos\alpha}{2}t,
\qquad 0<t<t_b=\sqrt{\frac{2L}{g\sin\alpha}}.}
\]
The upper endpoint is excluded: the target must not yet have reached $O$.

\textbf{b) Minimum without the deadline.} Write $c=g\cos\alpha/2$. Then
\[
v_0^2=\frac{L^2}{t^2}+c^2t^2
=\left(\frac Lt-ct\right)^2+2Lc\ge gL\cos\alpha.
\]
Equality requires $L/t=ct$, so
\[
\boxed{t_*=\sqrt{\frac{2L}{g\cos\alpha}},\qquad
v_{0,\min}=\sqrt{gL\cos\alpha}.}
\]
The two launch components are equal: the direction is $45^\circ$ above the ramp.

\textbf{c) Restore the deadline.} The optimum is possible precisely when
\[
t_*<t_b\quad\Longleftrightarrow\quad\tan\alpha<1
\quad\Longleftrightarrow\quad\alpha<45^\circ.
\]
For $\alpha\ge45^\circ$, the speed decreases throughout the allowed interval: $L/t-ct$ is positive and decreases there. Its limiting value occurs as $t\to t_b^-$:
\[
\boxed{\begin{aligned}
0<\alpha<45^\circ:&\quad v_{0,\min}=\sqrt{gL\cos\alpha}\quad\text{(attained)},\\[3pt]
45^\circ\le\alpha<90^\circ:&\quad v_0>\sqrt{\frac{gL}{2\sin\alpha}}\quad\text{(bound never attained)}.
\end{aligned}}
\]
At $45^\circ$ the optimum occurs exactly at the forbidden deadline. For steeper ramps it occurs too late. In both cases a successful launch can approach the bound as closely as desired, but there is no slowest one.

{\small\textbf{Flight check.} For any allowed meeting time $t$, at an earlier time $0<\tau<t$ the ball has
$n_B(\tau)=\tfrac12g\cos\alpha\,\tau(t-\tau)>0$.
It stays above the ramp until the meeting.}
\end{document}
